% Laget automatisk av elevark/lag_elevark.py fra sim/kuler-i-snorer.html. Endringer her blir overskrevet.
\documentclass[11pt]{article}
\usepackage{elevark}
\begin{document}
\elevarkhode{Fysikk 1 · Bevegelsesmengde og støt}{To kuler i snorer}{https://lektorbo.no/sim/kuler-i-snorer.html}{Simuleringen}

To kuler henger i hver sin snor på 1,0~m og så vidt berører hverandre. Kule A trekkes ut og slippes. Den treffer kule B i bunnen. Hvor høyt svinger kulene etter støtet?

\section*{Oppgaver}

Del forløpet i tre: A svinger ned, kulene støter, kulene svinger opp. Hvilken bevaringslov bruker du i hver del?

\begin{oppgaveliste}
\oppgave \fase{gjett}{Gjett}Kulene har like stor masse, og støtet er elastisk. A slippes fra 60°. Hva skjer med A og B etter støtet?
\svar
\oppgave \fase{test}{Test}Kjør simuleringen. Stemte tipset? Hvor høyt kommer B? Endre startvinkelen og se om det samme skjer.
\svar
\oppgave \fase{test}{Regn}A har massen 40~g og B 120~g. A slippes fra 60°, og støtet er elastisk. Regn ut farten til A rett før støtet og farten til begge kulene rett etter. Hvor stort blir utslaget til hver kule? Sjekk med simuleringen.
\svar
\oppgave \fase{gjett}{Gjett}Bytt til leirkuler med de samme massene. Hvor høyt svinger de sammen? Hvor stor del av den kinetiske energien forsvinner i støtet?
\svar
\oppgave \fase{forklar}{Forklar}Prøv en tung A mot en lett B, og så en lett A mot en tung B. Når spretter A tilbake, og når fortsetter den framover? Forklar ved hjelp av bevaring av bevegelsesmengde.
\svar
\oppgave \fase{forklar}{Forklar}La simuleringen gå en stund. Kulene støter igjen og igjen, men ikke alltid i bunnen. Hvorfor? Tips: Svingetiden til en pendel avhenger litt av utslaget.
\svar
\oppgave \fase{utfordring}{Utfordring}I et ekte forsøk med to kuler på 40~g og 120~g og startvinkel 60° svinger B ut til bare 26°. Bruk bevaring av bevegelsesmengde til å finne farten til A rett etter støtet. Hvor mye kinetisk energi gikk tapt i støtet, og hvor ble den av?
\svar
\end{oppgaveliste}

\end{document}
